% ============================================================
%  Chapter 09 — Just War Theory (During/After) · Quality of
%               Service Delivery · Code of Conduct & Ethics ·
%               Citizen Charter · Work Culture · Transparency &
%               Accountability · Social Audit · RTI · Probity in
%               Governance · Utilization of Public Funds
% ============================================================
\chapter{Service Delivery, Codes, Transparency, Accountability \& Probity}
\label{ch:service-delivery-codes-transparency-probity}

\section{Just War Theory --- During War and After War}
The Just War Theory provides an ethical framework that prescribes principles governing conduct before, during, and after war. The pre-war (jus ad bellum) principles were discussed previously. Below are the principles for during war (jus in bello) and after war (jus post bellum).

\subsection{During War (Jus in Bello) --- Principles}

\paragraph{Proportionality of Force.} During war, the force used must be proportional to the military objective. Disproportionate violence --- such as using excessive force against a minor target --- is ethically impermissible. Proportionality ensures that the destruction caused during conflict does not exceed what is necessary.

\begin{KeyConcept}
Proportionality of Violence: The degree of force used in war must be proportionate to the military objective being pursued. Excessive or indiscriminate force violates ethical norms of warfare.
\end{KeyConcept}

\paragraph{Distinction Between Civilian and Military Targets.} Combatants must distinguish between civilian and military targets. Deliberately targeting civilians or civilian infrastructure (hospitals, schools, places of worship) is a violation of the just war framework. Only legitimate military targets may be engaged.

\begin{CriticalPoint}
Combatants must distinguish between civilian and military targets. Intentional attacks on civilian populations constitute a war crime under international humanitarian law.
\end{CriticalPoint}

\paragraph{Humane Treatment of Prisoners of War and Wounded Soldiers.} Prisoners of war (POWs) and wounded soldiers should be treated humanely. Atrocities, torture, and inhumane treatment are strictly prohibited. This principle aligns with the Geneva Conventions.

\begin{KeyConcept}
Humane Treatment: Prisoners of war and wounded soldiers must be treated humanely. No atrocities or inhumane treatment should be inflicted upon them.
\end{KeyConcept}

\subsection{After War (Jus Post Bellum) --- Principles}

\paragraph{Holding Individuals Accountable for War Crimes.} After a war concludes, individuals responsible for war crimes must be identified and held accountable. International tribunals and courts of justice serve this function.

\paragraph{Establishing a Just and Lasting Peace.} The post-war phase must focus on establishing a just and lasting peace. This includes rebuilding societies, restoring institutions, ensuring justice for victims, and creating conditions that prevent future conflicts.

\begin{KeyConcept}
Post-war objectives include: (a) Accountability for war crimes; (b) Establishing a just and lasting peace; (c) Rebuilding societies and institutions.
\end{KeyConcept}

\textit{Note:} The before-war (jus ad bellum) principles are the most important from a decision-making perspective, as they deal with the ethical question of whether to go to war at all. During-war and after-war principles provide the ethical guardrails for conduct once war has begun.

\begin{QuickRevision}
During War: Proportionality of force; Distinction between civilian and military targets; Humane treatment of POWs and wounded soldiers.\\
After War: Holding individuals accountable for war crimes; Establishing a just and lasting peace; Rebuilding societies.\\
Before war (already discussed) involves the most critical decision-making about whether war is justified at all.
\end{QuickRevision}

\section{Quality of Service Delivery}

\subsection{Definition}
Quality of service delivery refers to the extent to which the services provided by an organization (especially government) meet or exceed the expected standards of citizens. When a service meets or exceeds expectations, it is considered qualitative. The topic is essentially about understanding the parameters or standards against which any service delivery can be evaluated.

\begin{KeyConcept}
Quality of Service Delivery: A service is considered qualitative when it meets or exceeds the expected standards. This topic involves understanding the parameters for evaluating any service delivery system.
\end{KeyConcept}

\subsection{Parameters of Quality of Service Delivery}

\paragraph{Availability.} Whether a service is available to all intended beneficiaries. If a section of society cannot access the service, it fails the availability test.

\begin{ExampleBox}
\begin{itemize}
\item Under the National Food Security Act, 5 kg of food grains were promised to certain populations through PDS. However, migrants could not avail these benefits because they were not at their registered locations. This was an availability problem.
\item The One Nation, One Ration Card scheme was introduced to address this availability gap, allowing migrants to access rations anywhere in the country.
\end{itemize}
\end{ExampleBox}

\paragraph{Accessibility.} Services must be accessible to all citizens, including persons with different abilities. Barriers to access must be identified and removed.

\begin{ExampleBox}
The Election Commission provides wheelchairs and ramps at polling stations for persons with disabilities, thereby increasing accessibility of the voting process.
\end{ExampleBox}

\paragraph{Responsiveness.} How quickly and effectively an organization responds to grievances is a key dimension. A responsive administration acknowledges complaints and resolves them promptly.

\begin{ExampleBox}
Many citizens who are dissatisfied with railway services tag railway officials on Twitter or social media, and their grievances are resolved through that channel --- demonstrating responsiveness through digital platforms.
\end{ExampleBox}

\paragraph{Reliability.} A service is reliable if citizens can depend on it consistently. If a service is unreliable, citizens must make alternative arrangements, which undermines trust.

\begin{ExampleBox}
\begin{itemize}
\item If one needs to book a Tatkal ticket for travel to Bihar during Chhath Puja, can one rely on the Tatkal system? If not, alternative arrangements are needed. Government reforms such as linking Tatkal booking with Aadhaar and fixing booking time windows were introduced to improve reliability.
\item Similarly, choosing an internet service provider depends on whose service is more reliable in a given area.
\end{itemize}
\end{ExampleBox}

\paragraph{Quality Per Se.} The inherent quality of the service itself matters greatly. If the quality of inputs or outputs is poor, the service fails regardless of other parameters.

\begin{ExampleBox}
\begin{itemize}
\item In PDS, it has been observed that good quality food grain is often sold in the open market while poor quality food grain is given to beneficiaries --- a direct quality failure.
\item Many people would not choose government schools or government hospitals for their families if given a choice, reflecting a perceived quality deficit in government service delivery.
\end{itemize}
\end{ExampleBox}

\paragraph{Effectiveness.} A service is effective if it achieves its intended objective. If a school does not deliver learning, or a hospital does not deliver treatment, the service is ineffective.

\begin{ExampleBox}
In many government schools, subject teachers (mathematics, physics, chemistry) are absent for long periods. This means learning is not happening, making the education service ineffective.
\end{ExampleBox}

\paragraph{Sustainability.} Services must be sustainable over time. Intermittent or temporary services are not truly qualitative.

\begin{ExampleBox}
If municipal water supply comes for only 9 months out of 12, it is not sustainable. Any government service must be consistently available over time.
\end{ExampleBox}

\paragraph{Economy, Efficiency, and Cost-Effectiveness.} If the cost of delivering a service is disproportionately high relative to the benefit (e.g., spending one rupee and only 16 paisa reaching the beneficiary), the service is not cost-effective and will not be sustainable in the long run.

\begin{ExampleBox}
Direct Benefit Transfer (DBT): Before DBT, benefits often did not reach beneficiaries' accounts due to corruption and leakages (efficiency problem). DBT improved both efficiency (reducing leakages) and effectiveness (ensuring the intended beneficiary actually receives the benefit).
\end{ExampleBox}

\begin{GovtPolicy}
Direct Benefit Transfer (DBT) improved both efficiency and effectiveness of government welfare delivery by transferring benefits directly to beneficiaries' bank accounts.
\end{GovtPolicy}

\paragraph{Ease and Convenience.} The ease with which a citizen can access a service is an important quality parameter. Reducing procedural complexity increases quality.

\begin{ExampleBox}
\begin{itemize}
\item The Jharkhand government explored home delivery of rations to improve ease of access.
\item UPSC introduced One Time Registration (OTR), eliminating the need for candidates to submit their information repeatedly across multiple examination cycles --- improving convenience.
\end{itemize}
\end{ExampleBox}

\paragraph{Timeliness.} Services must be delivered within a reasonable and expected timeframe. Delayed services lose their value and relevance.

\begin{ExampleBox}
Birth certificates, death certificates, and marriage certificates must be issued in a timely manner. A death certificate that is not issued promptly can cause serious legal and financial problems for the family of the deceased.
\end{ExampleBox}

\paragraph{Grievance Redressal.} A robust grievance redressal mechanism is essential. Citizens must have a channel to register complaints when services fail.

\begin{GovtPolicy}
CP-GRAM (Centralized Public Grievance Redress and Monitoring System): A government portal for grievance redressal related to central government departments. It has been integrated with AI (through a tool called `Samadhan Didi') for improved grievance resolution.
\end{GovtPolicy}

\paragraph{Citizen-Centricity.} Services must be designed around the needs of citizens. Private sector organizations like Domino's replace a pizza if the customer is dissatisfied --- a consumer-centric approach. Government services often lack this orientation.

\begin{CriticalPoint}
In the private sector, consumer-centricity drives accountability. In the government sector, citizens often have no alternative provider, making citizen-centric design even more important.
\end{CriticalPoint}

\begin{QuickRevision}
Parameters: Availability, Accessibility, Responsiveness, Reliability, Quality per se, Effectiveness, Sustainability, Economy/Efficiency/Cost-effectiveness, Ease/Convenience, Timeliness, Grievance Redressal, Citizen-centricity.\\
Key examples: One Nation One Ration Card (availability), Ramps at polling stations (accessibility), Social media grievance handling by railways (responsiveness), Tatkal reforms (reliability), PDS quality issues (quality per se), DBT (efficiency + effectiveness), OTR by UPSC (convenience), CP-GRAM with Samadhan Didi AI (grievance redressal).
\end{QuickRevision}

\subsection{Reasons for Poor Quality of Service Delivery}

\paragraph{Overburdening of the System / Population Pressure.} When the carrying capacity of public infrastructure is breached, service quality declines. Overburdened systems cannot maintain quality.

\begin{ExampleBox}
\begin{itemize}
\item Mumbai local trains are overcrowded due to massive population pressure, degrading the quality of transport.
\item Tourist places often face strain on sewer systems, drinking water supply, and waste management when tourist numbers exceed carrying capacity.
\end{itemize}
\end{ExampleBox}

\paragraph{Lack of Skilled Workforce.} When government employees, especially teachers and front-line workers, are not adequately trained, service delivery suffers. Many government teachers are not trained in sound pedagogy and rely on rote learning methods.

\begin{ExampleBox}
\begin{itemize}
\item Government teachers who rely on rote learning rather than conceptual teaching reduce the effectiveness of education delivery.
\item Skills India and Digital India initiatives aim to bridge the skill gap, but the training of government employees themselves remains a challenge.
\end{itemize}
\end{ExampleBox}

\paragraph{Corruption.} Corruption diverts resources meant for public services. Where corruption exists, poor service delivery will inevitably follow. Both supply-side (officials demanding bribes) and demand-side (citizens offering bribes) corruption contribute.

\paragraph{Bureaucratic Problems.} Several systemic issues within the bureaucracy lead to poor service delivery. These include colonial attitudes, top-down implementation, a culture of secrecy, red-tapism, resistance to change, and poor accountability mechanisms. These factors have been discussed in detail under bureaucratic challenges earlier in the syllabus.

\paragraph{Resource Constraints.} Inadequate investment in public services directly affects quality. Without sufficient financial resources, infrastructure remains poor and human resources are inadequate.

\begin{DataFact}
India's spending on healthcare and education as a percentage of GDP remains lower than that of many developed nations. For education, it is approximately 3\% of GDP, which is considered insufficient to provide quality services across the country.
\end{DataFact}

\paragraph{Digital Divide.} Many services are now delivered digitally, but significant portions of the population lack internet access, digital skills, or face vernacular language barriers. Where electricity infrastructure itself is absent, digital delivery is impossible.

\paragraph{Cultural and Societal Factors.} Even when the government provides good infrastructure, cultural factors can undermine service delivery. Poor civic sense, vandalism of public property, and normalization of wrong means all contribute to declining service quality.

\begin{ExampleBox}
\begin{itemize}
\item In trains, mugs in toilets are chained due to trust deficit --- people steal public property.
\item People throw stones at Vande Bharat trains, damaging public infrastructure.
\item Government can provide dustbins, but ultimately citizens must use them. Government can build toilets, but citizens must use them.
\item Corruption has a demand side too --- citizens themselves offer bribes, normalizing the practice.
\item Traffic law violations by citizens themselves degrade the quality of transport services.
\end{itemize}
\end{ExampleBox}

\begin{CriticalPoint}
Cultural and societal factors play a huge role: absence of good role models, normalization of wrong means, rising materialism, and poor civic sense all contribute to declining service quality despite government infrastructure.
\end{CriticalPoint}

\paragraph{Political Interference.} Political interference in tender processes, appointments, and implementation distorts service delivery. When contracts are given to political allies rather than the most competent bidder, quality suffers.

\begin{ExampleBox}
Crony capitalism, where contracts are awarded based on political connections rather than merit, is an example of how political interference degrades service delivery.
\end{ExampleBox}

\begin{QuickRevision}
Reasons for poor quality: Overburdening/population pressure, Lack of skilled workforce, Corruption, Bureaucratic problems, Resource constraints, Digital divide, Cultural/societal factors, Political interference.\\
Two broad categories: (1) Society-side factors (civic sense, cultural attitudes, demand-side corruption); (2) Governance/bureaucracy-side factors (colonial attitude, red-tapism, resource constraints, political interference).\\
Work culture topics (discussed below) add further bureaucratic challenges to this list.
\end{QuickRevision}

\section{Code of Conduct and Code of Ethics}

\subsection{The Three-Tiered Framework (Second ARC Recommendation)}
The Second Administrative Reforms Commission (ARC) recommended a three-tiered framework for guiding the behavior of civil servants. This framework is visualized as a triangle/pyramid:
\begin{itemize}
\item \textbf{Apex (Top)}: Statement of Values --- the broadest, most fundamental values (e.g., integrity, sustainability, justice).
\item \textbf{Middle}: Code of Ethics --- ethical standards and principles derived from those values. Broader than code of conduct.
\item \textbf{Base (Bottom)}: Code of Conduct --- specific do's and don'ts for particular situations. Most specific and detailed.
\end{itemize}

\begin{KeyConcept}
The Second ARC recommended three levels: (1) Statement of Values at the apex; (2) Code of Ethics derived from those values; (3) Code of Conduct as specific do's and don'ts. Each level becomes progressively more specific.
\end{KeyConcept}

\subsection{Understanding Code of Conduct}
Code of conduct refers to the specific do's and don'ts that clearly tell a civil servant what actions are permissible and which are not. It functions like a rulebook --- explicit, directive, and situation-specific.

\paragraph{Key Features of Code of Conduct.}
\begin{itemize}
\item Specifically tells you what to do and what not to do in clear terms.
\item In India, the Civil Services Conduct Rules serve as the code of conduct for civil servants.
\item Examples of specific rules: Civil servants cannot accept dowry; cannot criticize the government publicly; cannot accept gifts above a certain value; cannot accept lavish hospitality from friends; cannot help family members in official matters; cannot express personal opinions on political matters; must declare political activities.
\end{itemize}

\begin{ExampleBox}
\begin{itemize}
\item Civil servants are not supposed to take dowry.
\item Civil servants cannot criticize the government publicly.
\item There are specific limits on gifts that can be accepted.
\item Lavish hospitality from friends cannot be accepted.
\item Civil servants cannot help family members (bhaibanduon ki madad) in official matters.
\end{itemize}
\end{ExampleBox}

\paragraph{Code of Conduct and Conflict of Interest.} Code of conduct gives specific guidance on avoiding conflict of interest in different contexts:
\begin{itemize}
\item Recruitment/Interviews: If your relative is appearing for an interview, you must recuse yourself from the panel.
\item Tender Process: If you have a conflict of interest in a tender, you must record it and recuse yourself.
\item Judicial Proceedings: If a case involves someone you know, mere recording of the conflict of interest is sufficient in some contexts.
\end{itemize}
All these different do's and don'ts in different contexts are guided by the same code of ethics principle: there should not be any conflict of interest in decision-making.

\subsection{Understanding Code of Ethics}
Code of ethics defines the standards of behavior that should guide our actions. It is broader than code of conduct. While code of conduct tells you exactly what to do and what not to do, code of ethics provides the underlying ethical principles that guide behavior across all situations --- including unforeseen ones.

\begin{KeyConcept}
Code of Ethics: Standards of behavior that should guide our action. Broader than code of conduct. Provides ethical principles applicable across multiple situations, including those not specifically covered by rules.
\end{KeyConcept}

\subsection{Relationship Between the Three Tiers}
The three tiers are interconnected. A single value can generate multiple ethical standards, and each ethical standard can be applied as specific do's and don'ts across many situations.

\begin{ExampleBox}
\begin{itemize}
\item Value (Statement of Values): Integrity $\rightarrow$ Ethical Standard (Code of Ethics): Avoid conflict of interest $\rightarrow$ Code of Conduct: In recruitment --- recuse from panel if relative is a candidate; In tender process --- record conflict and recuse; In judicial proceedings --- record conflict.
\item Value: Sustainability $\rightarrow$ Code of Ethics: Ethical resource management, Wildlife conservation $\rightarrow$ Code of Conduct: No single-use plastic; E-waste buyback requirements; Project-specific conservation rules (e.g., Project Tiger, Project Elephant).
\end{itemize}
\end{ExampleBox}

\begin{ComparisonBox}
\textbf{Analogy with the Indian Constitution}:
\begin{itemize}
\item Preamble = Statement of Values (democratic, social, secular, republic, justice, equality, fraternity).
\item Directive Principles of State Policy (DPSPs) = Code of Ethics (broad ethical guidelines --- e.g., state should minimize inequality).
\item Fundamental Rights and specific legislation = Code of Conduct (specific enforceable rules --- e.g., taxation to reduce inequality).
\end{itemize}
\end{ComparisonBox}

\subsection{Code of Conduct is Application of Code of Ethics in Specific Situations}
Code of conduct is nothing but the application of code of ethics in specific situations that provides further clarity. If one simply states `avoid conflict of interest' (code of ethics), an individual can claim there is no conflict. To prevent such misuse of discretion, code of conduct provides specific rules.

\begin{CriticalPoint}
Code of conduct is an application of code of ethics in specific situations. Code of ethics is broader; code of conduct is more specific with clear do's and don'ts.
\end{CriticalPoint}

\subsection{Limitations of Code of Conduct}
The limitations of code of conduct overlap significantly with the limitations of laws, rules, and regulations (previously discussed in the `Law versus Ethics' topic).
\begin{itemize}
\item Cannot envisage all situations: Being a set of do's and don'ts, it cannot conceive all probable situations. New challenges (AI, blockchain, disaster management) may not be covered.
\item Discretion remains: Even with specific rules, some degree of discretion is unavoidable.
\item Implementation gaps: Like any rule, effective implementation is necessary. Without enforcement, rules remain on paper.
\item Lack of awareness: Many civil servants and citizens may not be fully aware of all conduct rules.
\item Vagueness in certain provisions: The rule that `a civil servant should not do anything that is unbecoming of a civil servant' is extremely broad and vague. `Unbecoming of a civil servant' can mean virtually anything.
\item Mixing of ethics and conduct: Rule 3 of conduct rules emphasizes integrity, duty, and ethical behavior --- these are essentially code of ethics principles placed within a code of conduct document, because India currently does not have a separate code of ethics.
\item Political interference: Political partisanship, improper appointments, and violations of conduct rules are common.
\item Weak enforcement and inadequate monitoring.
\end{itemize}

\begin{GovtPolicy}
Civil Services Conduct Rules: The rules that constitute the code of conduct for Indian civil servants. These are essentially `rules' --- and therefore share all the limitations of laws and regulations in ensuring ethical guidance.
\end{GovtPolicy}

\begin{CriticalPoint}
Key insight: When discussing limitations of code of conduct, the same arguments from `limitations of laws, rules, and regulations' (Law vs Ethics topic) apply directly, because code of conduct IS a set of rules. This overlap/congruency is a crucial theme for answer writing.
\end{CriticalPoint}

\subsection{Why Code of Ethics Is Important}
Given the limitations of code of conduct (which is essentially a set of rules), code of ethics becomes essential. The same reasoning that makes ethics important over mere laws applies here: laws and rules have limitations; ethical principles provide guidance where rules fall short, where discretion exists, and in unforeseen situations.

As of now, India does not have a formal code of ethics for civil servants. The Second ARC has recommended that there should be a statement of values and a code of ethics to complement the existing code of conduct.

\begin{tabularx}{\textwidth}{@{}>{\bfseries\color{Accent}}p{2.4cm} Y Y@{}}
\toprule
\rowcolor{AccentDark}\thead{Aspect} & \thead{Code of Conduct} & \thead{Code of Ethics} \\
\midrule
Nature & Specific do's and don'ts & Broad ethical standards/principles \\
Scope & Narrow --- covers specific situations & Broad --- guides behavior across all situations \\
Flexibility & Rigid --- prescriptive rules & Flexible --- allows judgment in unforeseen situations \\
Example & Cannot accept gifts above \rupee X; must recuse from panel if relative is candidate & Avoid conflict of interest; uphold integrity \\
Indian Context & Civil Services Conduct Rules & No formal code of ethics exists yet (recommended by 2nd ARC) \\
Limitation & Cannot cover all situations; vague in places & May be too broad; depends on individual interpretation \\
\bottomrule
\end{tabularx}

\begin{PYQLink}
UPSC has asked about the difference between code of conduct and code of ethics multiple times. The three-tiered approach (Statement of Values $\rightarrow$ Code of Ethics $\rightarrow$ Code of Conduct) recommended by the Second ARC is a frequently tested concept.
\end{PYQLink}

\begin{QuickRevision}
Three tiers (Second ARC): Statement of Values (apex) $\rightarrow$ Code of Ethics (middle) $\rightarrow$ Code of Conduct (base/most specific).\\
Code of Conduct = Specific do's and don'ts (Civil Services Conduct Rules in India). Code of Ethics = Broader ethical standards that guide action across all situations.\\
Code of Conduct is an application of Code of Ethics in specific situations.\\
Limitations of Code of Conduct = Limitations of laws and regulations (Law vs Ethics topic applies directly).\\
India lacks a formal Code of Ethics --- Second ARC recommends creating one.\\
Key vagueness: `Unbecoming of a civil servant' --- too broad, defeats the purpose of specificity.
\end{QuickRevision}

\section{Citizen Charter}

\subsection{Background and Context}
Traditionally, bureaucracy operated under a closed-system approach where citizens were seen merely as passive recipients of goods and services --- not as active participants in governance. The administration implemented policies made by the political executive without evaluating whether those services were actually relevant to or beneficial for citizens. This was not citizen-centric administration.

Under the New Public Management (NPM) approach, originating in the United Kingdom, the idea emerged that administration should be made citizen-centric --- borrowing the efficiency and accountability mechanisms of the private sector. In the private sector, consumers are at the center of service delivery; organizations are responsive because they know that ignoring consumers will have consequences.

In government services, however, there was no such orientation. Whether a government school teacher's teaching was effective, or a government doctor's treatment was working, was often not the administration's concern. The perspective that `all services exist to serve citizens' was missing.

\begin{KeyConcept}
Citizen Charter emerged from the New Public Management approach, which sought to make government administration citizen-centric --- just as private sector organizations are consumer-centric. It is a partnership document between an organization and its citizens.
\end{KeyConcept}

\subsection{Definition and Components}
A citizen charter is an agreement (a partnership) between an organization and citizens, in which the organization specifies:
\begin{itemize}
\item Quantity and quality of services it will provide.
\item Expected timelines for service delivery.
\item Grievance redressal mechanism if services are not delivered as promised.
\end{itemize}
Essentially, the organization enters into a partnership with citizens, placing them at the center of administration and promising: `This is what you can expect from me, and if I fail, here is how you can hold me accountable.'

\subsection{Case Studies}

\paragraph{NTPC Citizen Charter.} NTPC's citizen charter includes:
\begin{itemize}
\item Vision: To be the world's leading power company.
\item Mission: To provide reliable power and related solutions.
\item Core Values: Integrity, consumer focus, etc.
\item Services to various stakeholders: Government (compliance, reports), Shareholders (dividends, annual reports), Consumers/State governments (supply of power).
\item Expectations from stakeholders: Government should provide timely clearances; contractors/vendors must adhere to safety guidelines and provide quality services; optimum utilization of resources.
\item Grievance redressal: Executive Director designated as head for grievance redressal; charter is reviewed periodically.
\end{itemize}

\begin{KeyConcept}
NTPC Citizen Charter provides for grievance redressal mechanism, expectations from stakeholders, and services provided to stakeholders. It demonstrates the partnership model of citizen charter.
\end{KeyConcept}

\paragraph{RBI Citizen Charter.} The RBI citizen charter provides timelines for various services such as loan processing, localized branch services, and other banking-related services.

\paragraph{Passport Seva Kendra.} The passport citizen charter specifies timelines for passport processing --- telling citizens exactly how many days different types of passport applications will take.

\paragraph{Hospitals (e.g., Gangaram).} Hospital citizen charters detail patients' rights (right to know diagnosis, available treatment options) and patients' responsibilities (paying fees on time, treating doctors respectfully, not engaging in violence).

\subsection{Purpose and Significance}
\begin{itemize}
\item Ensures administration is accountable to citizens.
\item Ensures administration truly serves the needs of citizens.
\item Enables citizen participation in governance.
\item Shifts the attitude from rule-implementation to citizen-centric service delivery.
\end{itemize}

\subsection{Limitations of Citizen Charter}
\begin{itemize}
\item Lack of Legal Enforcement: Citizen charters do not have legal backing. There is no law mandating them, leading to poor accountability.
\item Non-existent Citizen Charters: Many organizations simply do not have citizen charters at all.
\item One-Size-Fits-All Approach: When a ministry adopts a single citizen charter, all departments under it follow the same template --- even though school education and higher education, for example, deliver very different services.
\item Lack of Citizens' Participation While Formulating: If citizens are not consulted during formulation, the charter will not capture their actual expectations.
\item Linguistic Barriers: Many citizen charters are written only in English. Those who most need the charter often do not understand English. Charters are often not available in vernacular languages.
\item Absence of Grievance Redressal Mechanism: Many charters make promises but do not provide a recourse mechanism if promises are not fulfilled.
\item Lack of Capacity Building: Before making promises (e.g., service in 15 days), the organization must have adequate capacity (staff, infrastructure). Often, promises are made without corresponding capacity building.
\item Bureaucratic Mindset / Colonial Attitude: The colonial, top-down, self-centered approach and culture of secrecy in the bureaucracy means they resist the citizen-centric orientation that citizen charters demand.
\item Inertia to Change and Red-Tapism.
\item Culture of Secrecy in the bureaucracy.
\item Lack of Awareness Among Public: Citizens are often not aware that citizen charters exist or what rights they confer.
\end{itemize}

\subsection{Measures to Improve (Second ARC Recommendations)}
\begin{itemize}
\item Legal framework for enforcement of citizen charters.
\item Consultation and participation of citizens during formulation.
\item Standardization of charter formats while allowing adaptation by individual departments.
\item Monitoring and evaluation mechanisms.
\item Capacity building and training for staff.
\item Effective grievance redressal mechanisms.
\item Charters in vernacular languages.
\end{itemize}

\begin{GovtPolicy}
Sevottam Model: A model for excellence in service delivery that provides a systematic approach to improve citizen charters. It is a good case study for exam answers.
\end{GovtPolicy}

\begin{QuickRevision}
Citizen charter = agreement between organization and citizens specifying services, timelines, and grievance mechanisms.\\
Background: Emerged from New Public Management; borrowed consumer-centric approach from private sector.\\
Case studies: NTPC (stakeholder services + grievance redressal), RBI (service timelines), Passport Seva Kendra (processing timelines), Hospitals (patient rights and responsibilities).\\
Key limitations: No legal backing, non-existent in many organizations, one-size-fits-all, no citizen participation in formulation, linguistic barriers, absent grievance redressal, bureaucratic mindset.\\
Measures: Legal framework, citizen consultation, standardization with flexibility, capacity building, grievance redressal, vernacular languages.\\
Sevottam Model --- a model for excellence in service delivery (good case study for answers).
\end{QuickRevision}

\section{Work Culture}

\subsection{Definition}
Work culture refers to the shared values, attitudes, behaviors, and practices that define how work is conducted within an organization. It is not determined by a single person's attitude --- it is defined by the collective values shared by employees. However, a single person (especially a leader) can change the work culture of an entire organization.

\begin{KeyConcept}
Work Culture = Shared values, attitudes, and behaviors that collectively define how an organization operates. It is collective, not individual --- though a single leader can transform it.
\end{KeyConcept}

\subsection{Dimensions of Work Culture}
\begin{itemize}
\item Leadership style: Authoritarian (top-down) vs. participatory (collaborative decision-making).
\item Communication: Open and transparent vs. closed and hierarchical.
\item Technology and innovation: Adoption of new technologies and openness to innovative approaches.
\item Decision-making: Centralized vs. decentralized; fast vs. slow.
\item Motivation and morale: How employees are motivated; collective morale levels.
\item Accountability culture: Whether accountability is enforced or evaded.
\item Transparency: Information sharing within the organization.
\item Division of labor: How tasks are distributed.
\item Hierarchy: Flat vs. deep hierarchical structures (related to Max Weber's bureaucratic model).
\item Recognition and appreciation: Whether good work is recognized and rewarded.
\end{itemize}

\subsection{Positive Work Culture --- Characteristics and Examples}
A healthy/positive work culture is characterized by clear and open communication, trust and respect among employees, a sense of belonging, opportunities for growth and development, a sense of purpose, and recognition and appreciation for good work.

\begin{ExampleBox}
\begin{itemize}
\item Delhi Metro Rail Corporation (DMRC): Under E.\ Sreedharan's leadership, DMRC developed a work culture of timely project delivery and punctuality. Despite government departments being notorious for delays, Sreedharan changed the culture so that all employees believed projects must be delivered on time. This is a classic example of positive work culture --- shared values of efficiency and punctuality.
\item Google: Trusts employees; provides nap pods, flexible working; focuses on targets rather than micromanagement; promotes innovation; rewards employees for sincere effort even when mistakes happen. Google also provides three months of paternity leave.
\item Tata Group: Known for ethical corporate governance, respect for employees and consumers, strong ethical values as the foundation of its business culture.
\end{itemize}
\end{ExampleBox}

\subsection{Negative Work Culture --- Characteristics and Examples}
A negative/toxic work culture is characterized by poor communication, lack of trust, high pressure without support, absence of work-life balance, corruption, and lack of recognition.

\begin{ExampleBox}
\begin{itemize}
\item Satyam Scam: A case study of negative work culture where lack of objectivity, poor corporate governance, and fraud became part of the organizational culture.
\item Glass Ceiling: Discrimination against women in promotions and career advancement is an example of negative work culture, both in corporate and government settings.
\item CRPF: High attrition rate and suicides among personnel due to work pressure --- an example of negative work culture and poor work-life balance.
\item Police: Working under high pressure, not getting adequate leave --- a negative example of work-life balance and work culture.
\item Work from Home concerns: Erosion of work-life boundaries; employees expected to respond at all hours; meetings scheduled late at night --- privacy is destroyed and mental health is impacted. This highlights the need for a `Right to Disconnect.'
\end{itemize}
\end{ExampleBox}

\subsection{Determinants of Work Culture}
In the government context, both corporate governance principles (private sector) and foundation values for civil services (government sector) influence work culture. The factors that influence work culture in the government include those discussed under bureaucratic challenges and under value inculcation in society.

\begin{GovtPolicy}
Right to Disconnect: The concept that employees should have the right to not respond to work communications outside of working hours. This is relevant in the context of work-from-home culture and its impact on mental health and work-life balance.
\end{GovtPolicy}

\begin{QuickRevision}
Work culture = shared values defining how work is done in an organization.\\
Dimensions: Leadership style, communication, technology, decision-making, motivation, accountability, transparency, hierarchy, recognition.\\
Positive examples: DMRC (timely delivery, E.\ Sreedharan), Google (trust, flexibility, innovation), Tata (ethical governance).\\
Negative examples: Satyam Scam (fraud culture), Glass ceiling (discrimination), CRPF (work pressure/suicides), Police (poor work-life balance), WFH without Right to Disconnect.\\
Work culture can be evaluated positively or negatively across any parameter; good or bad case studies can be cited.
\end{QuickRevision}

\section{Transparency and Accountability}

\subsection{Transparency}

\paragraph{Definition.} Transparency is about making information available to the people --- timely, reliably, and in an accessible manner --- so that it can be used for scrutiny and in the interest of the public. When citizens have information, they are empowered and can hold the government accountable.

\begin{KeyConcept}
Transparency: Making information available to the people in a timely, reliable, and accessible manner so that it serves as a tool for scrutiny and public accountability.
\end{KeyConcept}

\paragraph{Types of Transparency.} \textbf{Event Transparency}: Transparency about inputs, outputs, and outcomes. For example: a road was constructed using \rupee 1 crore, which produced 5 km connecting Village A and Village B. This tells you what was spent and what was achieved.

\textbf{Process Transparency}: Transparency about the procedures and processes followed. For example: how was the particular contractor (ABC Limited) selected? Was there a fair bidding process? If ABC Limited is the minister's son's company, there is a problem. Process transparency ensures citizens know how decisions were made --- not just what the outcomes were.

\begin{ExampleBox}
UPSC mostly follows event transparency --- it publishes results, marks, and answer keys (input, output, outcome). But exactly how copies are being checked, what methods are used for evaluation --- that procedural information is not shared. That would be process transparency.
\end{ExampleBox}

\begin{CriticalPoint}
Transparency should not merely provide token information. The information must be meaningful enough for citizens to use it for accountability. Both event transparency AND process transparency are needed for genuine empowerment.
\end{CriticalPoint}

\subsection{Accountability}

\paragraph{Definition.} Accountability is the obligation --- especially of those in decision-making positions --- to report, explain, and be answerable for their actions and inactions, and to face the resulting consequences. It involves an associated punishment or sanction mechanism for non-fulfillment of obligations.

\begin{KeyConcept}
Accountability: The obligation of decision-makers to explain their actions/inactions and face consequences. It includes: (1) Obligation to explain; (2) Answerability; (3) Facing consequences (punishment/sanctions).
\end{KeyConcept}

\paragraph{Accountability vs. Responsibility.}

\begin{tabularx}{\textwidth}{@{}>{\bfseries\color{Accent}}p{2.2cm} Y Y@{}}
\toprule
\rowcolor{AccentDark}\thead{Aspect} & \thead{Accountability} & \thead{Responsibility} \\
\midrule
Timing & Usually after the task is completed & During the course of the task \\
Assignment & Fixed on specific person(s) --- assigned authoritatively & Can be shared; can be assumed voluntarily \\
Nature & Cannot be taken voluntarily --- must be assigned with prior knowledge of obligations & Can be taken upon oneself (e.g., collective responsibility for environment) \\
Example & DM/SP are held accountable after an incident --- they knew beforehand they would answer & Keeping classroom clean --- collective responsibility assumed by all \\
Scope & Specific to assigned duties and expectations & Broader --- can extend to moral and social duties \\
\bottomrule
\end{tabularx}

\subsection{Traditional Forms of Accountability}
Traditional accountability mechanisms include:
\begin{itemize}
\item \textbf{Legal Accountability}: Laws such as the Prevention of Corruption Act that punish wrongdoers.
\item \textbf{Political Accountability}: Opposition in Parliament holding the government accountable through questions, debates, and no-confidence motions.
\item \textbf{Administrative Accountability}: Hierarchy (Max Weber's model) --- superiors holding subordinates accountable. E.g., DC holds DSP accountable, DSP holds Inspector accountable.
\item \textbf{Fiscal Accountability}: Institutions like the Comptroller and Auditor General (CAG) that audit government expenditure.
\item \textbf{Judicial Accountability}: Courts pronouncing judgments and fixing punishment.
\item \textbf{Separation of Powers}: Each branch of government checking the others.
\end{itemize}

\subsection{Limitations of Traditional Accountability}
\begin{itemize}
\item Inward-looking/Institutional-centric: Traditional mechanisms evaluate policies from the administration's perspective, not from the citizens' perspective. Closed system approach.
\item Slow decision-making: Judiciary is slow; cases take years to resolve.
\item Reactive rather than proactive: Issues are addressed only after they become problems (e.g., NEET leak controversy in Parliament). CAG does post-mortem auditing.
\item Focus on legality, not outcomes: They check whether procedures were legally followed, not whether citizens actually benefited.
\item Not people-centric: Citizens' perspectives and actual impact on their lives are often not evaluated.
\end{itemize}

\begin{CriticalPoint}
Traditional accountability mechanisms are institutional-centric, slow, reactive, and focused on legality rather than outcomes. This creates a gap that social accountability mechanisms must fill.
\end{CriticalPoint}

\subsection{Social Accountability Mechanisms}
Given the limitations of traditional accountability, social accountability mechanisms have emerged to complement them. Social accountability involves citizens actively participating in holding government accountable.
\begin{itemize}
\item Citizen Charter: Empowers citizens to question service delivery.
\item RTI (Right to Information): Allows citizens to access government information.
\item Social Audit: Citizens collectively monitor and evaluate government schemes.
\item Social Media: Debates, discussions, and protests on social media platforms.
\item Citizens' Report Card: Citizens rate government services.
\item Public Interest Litigation (PIL): Citizens approach courts on matters of public interest.
\item Jan Sunwai (Public Hearings): Citizens participate in hearing processes.
\item Protest movements and Janta Darbar.
\item Participatory budgeting.
\item NGO performance reports (e.g., ASER Report on learning levels).
\end{itemize}

\begin{QuickRevision}
Transparency: Making information available --- both event transparency (inputs/outputs/outcomes) and process transparency (procedures/decisions).\\
Accountability: Obligation to explain, be answerable, and face consequences. Different from responsibility (which can be shared and assumed).\\
Traditional accountability: Legal, political, administrative, fiscal, judicial. Limitations: institutional-centric, slow, reactive, legality-focused.\\
Social accountability: Citizen charter, RTI, social audit, social media, PIL, citizens' report card, jan sunwai, participatory budgeting, NGO reports (ASER).
\end{QuickRevision}

\section{Social Audit}

\subsection{Definition}
Social audit is a participatory process through which citizens collectively monitor, verify, and evaluate the implementation of government schemes by comparing official records with the ground reality, in order to ensure accountability, transparency, and participation.

\begin{KeyConcept}
Social Audit: A participatory process in which citizens collectively monitor and evaluate government schemes by comparing official records with ground reality --- ensuring accountability, transparency, and citizen participation.
\end{KeyConcept}

\subsection{Why Social Audit Is Needed}
Traditional accountability mechanisms only check whether procedures were legally followed and money was legally spent. They do not assess whether citizens actually benefited from the expenditure. Social audit fills this gap by evaluating schemes from the beneficiaries' perspective.

\begin{ExampleBox}
Under MNAREGA: Traditional accountability would check whether \rupee 20{,}000 crore was legally disbursed, budgeted, and procedurally correct. But social audit checks whether the beneficiaries actually received benefits, whether the work was actually done on the ground, and whether the listed beneficiaries actually exist.
\end{ExampleBox}

\subsection{How Social Audit Works (MNAREGA Example)}
\begin{itemize}
\item Gram Sabha meetings are convened with participation of all villagers.
\item The list of MNAREGA beneficiaries is examined by the community.
\item Villagers identify ghost beneficiaries --- e.g., names like Amitabh Bachchan, Shahrukh Khan, Sunny Leone found in beneficiary lists, which are obviously fake identities used for embezzlement.
\item Physical verification of work: If a pond was supposed to be built, villagers can see whether it actually exists.
\item Comparison of official records with ground reality exposes corruption, fake muster rolls, and diversion of funds.
\end{itemize}

\subsection{Case Study: Meghalaya}
\begin{GovtPolicy}
Meghalaya Society for Social Audit and Transparency (MSSAT): Meghalaya has enacted a law making social audit mandatory for a number of schemes and departments. This is a model case study for social audit implementation.
\end{GovtPolicy}

\begin{ExampleBox}
\textbf{Case examples from Meghalaya's social audit}
\begin{itemize}
\item Old Age Pension Case: A beneficiary's pension under the National Social Assistance Programme was stopped because his name was erroneously removed from the database --- officially recorded as deceased. Social audit revealed the person was alive, and pension was resumed. (A similar case happened in Haryana where a person took out a procession with a banner saying `Tumhara Phuppa Zinda Hai' --- your uncle is alive.)
\item Non-Functioning School Case: A school had not been functioning for three years. Midday meals were not provided. Parents had admitted children to other schools. Social audit revealed this, and the education department decided to recover the misused funds (\rupee 84{,}000 reimbursed to the department). This demonstrates holding administration accountable for non-functioning of midday meals.
\end{itemize}
\end{ExampleBox}

\subsection{Benefits of Social Audit}
\begin{itemize}
\item Enhances transparency in government operations.
\item Strengthens accountability of officials and institutions.
\item Promotes citizen participation in governance.
\item Detects corruption and leakages.
\item Improves service delivery quality.
\item Empowers beneficiaries to demand their rights.
\item Builds public trust in government institutions.
\item Provides continuous feedback for course correction.
\item Strengthens grassroots democracy through Gram Sabha participation.
\item Promotes ethical governance.
\item Improves government legitimacy.
\end{itemize}
Additionally, utilitarian arguments can be added --- social audit enhances the greatest good for the greatest number by ensuring resources reach actual beneficiaries.

\subsection{Limitations of Social Audit}
Many limitations overlap with those of citizen charter, since both lack comprehensive legal enforcement.
\begin{itemize}
\item Low public awareness about social audit processes.
\item Political interference --- local leaders (gram panchayat) may resist social audit to protect corruption.
\item Bureaucratic resistance --- attitude of secrecy, reluctance to share information.
\item Poor quality or inaccessible records --- officials may not maintain proper records or may refuse to share them.
\item Inadequate institutional capacity and resources.
\item Weak follow-up mechanism --- no enforcement of audit findings.
\item Lack of enforcement power --- not legally mandated for most schemes (MNAREGA is an exception).
\item Elite capture --- powerful local elites manipulate the process to protect their interests.
\item Limited capacity of citizens to analyze government documents.
\end{itemize}

\begin{CriticalPoint}
Social audit is highly probable for UPSC Mains this year. Along with Citizen Charter and RTI, it is one of the three most likely topics for a question in GS Paper IV.
\end{CriticalPoint}

\begin{QuickRevision}
Social audit = participatory process; citizens compare official records with ground reality.\\
Needed because traditional accountability is institutional-centric and does not check citizen-level outcomes.\\
MNAREGA: Gram Sabha examines beneficiary lists, physically verifies work, identifies ghost beneficiaries.\\
Case study: Meghalaya Society for Social Audit and Transparency (MSSAT) --- legally mandated social audit.\\
Benefits: Transparency, accountability, participation, corruption detection, empowerment, public trust, grassroots democracy, ethical governance.\\
Limitations: No legal backing (except MNAREGA), low awareness, political interference, bureaucratic resistance, weak follow-up, elite capture.
\end{QuickRevision}

\section{Right to Information (RTI)}

\subsection{Overview}
RTI gives every citizen the right to access information held by public authorities. It is a key social accountability tool that empowers citizens to hold government accountable.
\begin{itemize}
\item Coverage: Applies to public institutions. Every department must appoint a Public Information Officer (PIO) to receive and process RTI applications.
\item Timeline: Usually 30 days for response. If life/liberty is threatened, 48 hours.
\item Proactive Disclosure: The Act mandates that organizations should proactively disclose certain information (e.g., UPSC publishes past year papers, declares answer keys).
\item No Reason Required: The applicant does not need to state why they are seeking the information. This prevents intimidation --- if reasons were required, the information could be used against the applicant.
\item Appeal Mechanism: First appeal within the department (to a senior officer). Second appeal to the Central/State Information Commission (CIC/SIC).
\item Exemptions: Certain information related to strategic interests, trade secrets, sovereignty, and national security is exempted.
\end{itemize}

\subsection{RTI as a Social Accountability Tool}
RTI enables citizens to:
\begin{itemize}
\item Verify tender processes --- check if eligibility criteria were changed to favor a particular party.
\item Access beneficiary lists --- verify if benefits reached actual beneficiaries or were diverted.
\item Check minutes of meetings --- understand how decisions were made.
\item Expose corruption and scams (e.g., the Adarsh Society scam was exposed through RTI).
\item Improve service delivery by creating transparency pressure.
\item Identify ghost beneficiaries in welfare schemes.
\end{itemize}

\subsection{Benefits of RTI}
\begin{itemize}
\item Increases transparency in government operations.
\item Strengthens accountability of officials.
\item Promotes citizen participation in governance.
\item Helps detect corruption and leakages.
\item Improves service delivery through transparency pressure.
\end{itemize}
\textit{Note:} The arguments for benefits of RTI are similar to those for social audit and citizen charter --- they all serve as social accountability mechanisms.

\subsection{Limitations and Challenges of RTI Implementation}

\paragraph{Infrastructural Issues.}
\begin{itemize}
\item Poor record management: Officers do not properly maintain or categorize records. When RTI is filed, the entire file has to be searched, consuming resources and time.
\item Lack of digitization of records.
\item Poor infrastructure for processing RTI applications in many departments.
\end{itemize}

\paragraph{Bureaucratic Attitude.}
\begin{itemize}
\item Officers are trained in secrecy --- their default attitude is to share minimum information.
\item Bureaucratic resistance to disclosure of information.
\item Fear of scrutiny deters officials from frank decision-making --- knowing that deliberations may be exposed via RTI, officers may avoid giving honest opinions in files, impacting administrative efficiency and innovation.
\end{itemize}

\paragraph{Capacity and Institutional Issues.}
\begin{itemize}
\item Frequent transfers of PIOs: PIOs need time to develop the skill of deciding what information to divulge and what to withhold. Before they gain adequate experience, they are transferred --- leading to capacity gaps.
\item CIC and SIC composition: Central and State Information Commissions are often dominated by retired bureaucrats who have been historically trained in protecting information, not disclosing it. This creates a `parking lot' mentality where these positions are used as post-retirement sinecures.
\item Lack of independence: The government can influence salaries and terms of information commissioners, compromising their independence. Original provisions of the RTI Act did not allow government to change these terms, but subsequent amendments introduced this power.
\item Overburdened Information Commissions: Vacancies in Information Commissioner positions are not filled timely, creating backlogs.
\end{itemize}

\paragraph{Enforcement Challenges.}
\begin{itemize}
\item Penalties are not strictly enforced --- poor deterrence for non-compliance.
\item High rejection rate in certain departments with frivolous reasons for rejection.
\item Rejection using broad exemption clauses like `wider public interest.'
\end{itemize}

\paragraph{Citizen-Side Challenges.}
\begin{itemize}
\item Complex application process and financial burden of filing RTI for common citizens.
\item Lack of awareness among citizens about RTI rights and procedures.
\item Misuse of RTI: Some people file frivolous RTI applications (e.g., asking how many samosas were served in a meeting) or use RTI for personal vendettas. This diverts resources, reduces efficiency, and undermines RTI's spirit.
\item RTI used as harassment tool against officers --- filing RTI to access personal opinions of officials on file, which chills honest deliberation.
\end{itemize}

\paragraph{Safety Concerns.} Threat to life of RTI activists: Many RTI activists have been attacked or killed for exposing corruption. Lack of protection for whistleblowers compounds this problem.

\paragraph{Impact on Administrative Efficiency.} When RTI is misused, officers spend excessive time responding to frivolous applications. This diverts resources from actual work, reduces efficiency, and makes officers risk-averse in their decision-making. The Economic Survey has also noted that RTI's impact on decision-making deliberations discourages frank expression of opinions, which affects policy quality.

\begin{CriticalPoint}
The impact of RTI misuse on administrative efficiency is an important point for answer writing: officers avoid giving honest opinions knowing deliberations may be disclosed, which impacts decision-making quality and innovation.
\end{CriticalPoint}

\begin{QuickRevision}
RTI = right to access information from public authorities. 30-day timeline; no reason required; appeal to CIC/SIC.\\
Social accountability tool: verifies tender processes, exposes corruption (Adarsh scam), identifies ghost beneficiaries.\\
Limitations: Poor records, bureaucratic secrecy, PIO transfers, CIC/SIC composition (retired bureaucrats), lack of independence, penalties not enforced, high rejection rates.\\
Citizen-side: Complex process, misuse (frivolous/harassment RTIs), low awareness.\\
Safety: Threats to RTI activists, lack of whistleblower protection.\\
Impact: Chills frank decision-making, diverts resources, reduces efficiency.
\end{QuickRevision}

\section{Probity in Governance}

\subsection{Probity vs.\ Integrity}

\paragraph{Probity.} Probity is the term used in the context of organizations and professions. It is about maximizing ethical principles in organizational conduct --- going beyond personal honesty to actively create systems, checks, and balances that minimize corruption and maximize ethical behavior in the organization.

\begin{KeyConcept}
Probity: Maximizing ethical principles in organizational/professional life. It includes leadership, creating checks and balances, and institutional mechanisms to minimize corruption. Probity mechanisms include: transparency improvements, accountability systems, RTI, social audit, DBT, and all systems designed to minimize corruption in governance.
\end{KeyConcept}

\paragraph{Integrity.} Integrity is about consistently following your principles and doing the right thing even when no one is watching. It applies to both personal and professional life. It is about being true to one's values across all contexts.

\begin{tabularx}{\textwidth}{@{}>{\bfseries\color{Accent}}p{2cm} Y Y@{}}
\toprule
\rowcolor{AccentDark}\thead{Aspect} & \thead{Probity} & \thead{Integrity} \\
\midrule
Context & Organizational and professional & Personal and professional \\
Focus & Maximizing ethical principles; creating systems & Consistency of personal values and actions \\
Scope & Institutional --- checks, balances, mechanisms & Individual --- character, honesty \\
Example & DBT, social audit, transparency mechanisms, anti-corruption institutions & A person not cheating even when no one is watching \\
Usage & `Lack of probity in governance' --- refers to organizational/systemic failure & `Lack of integrity' --- refers to personal moral failure (e.g., disloyalty) \\
\bottomrule
\end{tabularx}

\subsection{Reasons for Declining Probity}
The reasons for declining probity can be categorized into several dimensions:

\paragraph{Political Dimension.}
\begin{itemize}
\item Political patronage and protection of corrupt officials.
\item Political interference in administration.
\item Lack of political will to implement probity mechanisms.
\end{itemize}

\paragraph{Administrative/Bureaucratic Dimension.}
\begin{itemize}
\item Poor work culture in government organizations.
\item Weak traditional accountability mechanisms.
\item Red-tapism and procedural bottlenecks.
\item Colonial attitudes and culture of secrecy.
\end{itemize}

\paragraph{Societal Dimension.}
\begin{itemize}
\item Normalization of corruption in society.
\item Poor role models --- corruption glorified in movies and popular culture.
\item Rising materialism and `ends justify means' approach.
\item Success measured only by wealth, regardless of how it was earned.
\item Demand-side corruption --- citizens themselves offer bribes.
\item Poor civic sense and breaking of rules by citizens themselves.
\end{itemize}

\paragraph{Moral/Individual Dimension.}
\begin{itemize}
\item Declining personal values and ethics.
\item Rising greed and selfishness.
\item Lack of moral courage to resist corruption.
\end{itemize}

\paragraph{Historical/Cultural Dimension.} Bakshish system from the Mughal era --- the cultural expectation of receiving extra rewards for performing one's duty. This has evolved into a form of bribery where officials expect `something extra' from citizens for doing their regular work (e.g., patwaris expecting payment for routine land records work).

\paragraph{Deterrence Dimension.}
\begin{itemize}
\item Poor deterrence --- insufficient punishment for corruption.
\item Slow judicial process means corrupt officials are rarely punished quickly.
\item Society does not create social pressure against corruption.
\end{itemize}

\subsection{Philosophical Basis of Probity in Governance}
This topic connects to all the philosophical perspectives on governance that have been studied under various thinkers:
\begin{itemize}
\item \textbf{Utilitarianism} (Bentham, Mill): Governance should maximize good for the maximum number of people. Lincoln's statement about government `of the people, by the people, for the people' reflects this.
\item \textbf{Deontology} (Kant): Human beings should be treated never merely as a means --- dignity and respect in governance.
\item \textbf{Virtue Ethics} (Aristotle): Golden mean; those who rule should be persons of wisdom and virtue (Plato's Philosopher King).
\item \textbf{Max Weber}: Rational-legal authority with written rules, hierarchy, and division of labor for accountability.
\item \textbf{Social Contract} (John Locke): Government authority is limited; derives legitimacy from consent of the governed.
\item \textbf{John Rawls}: Just institutions; veil of ignorance; inequality should favor the weaker sections.
\item \textbf{Machiavelli}: Politics is about power (his philosophical perspective, though contrasting with ethical governance ideals).
\item \textbf{Gandhi}: `Politics without principle is a sin'; Gandhian Talisman --- test every policy by its impact on the poorest.
\item \textbf{Bhagavad Gita and Mahabharata}: Philosophical foundations rooted in Indian traditions about dharma and governance.
\end{itemize}

\paragraph{Kautilya (Arthashastra).} Kautilya's philosophy on governance contains crucial insights on probity:
\begin{itemize}
\item `In the happiness of his subjects lies the king's happiness; in their welfare, his welfare.' --- Administration must be citizen-centric, not self-serving.
\item Corruption is inevitable --- `Just like you cannot stop fish from drinking water, you cannot stop corruption.' Therefore, the focus should be on probity mechanisms to minimize corruption.
\item Recommended a network of spies (guptachars) for detection and hefty punishment for deterrence.
\item Sam, Dam, Dand, Bhed --- different strategies for governance and maintaining order.
\end{itemize}

\begin{KeyConcept}
Kautilya's key insight: Since corruption is inevitable, the focus must be on creating robust probity mechanisms --- detection systems, checks and balances, and deterrent punishments. This is the philosophical basis for all modern anti-corruption institutions.
\end{KeyConcept}

\begin{QuickRevision}
Probity = maximizing ethical principles in organizational context (not just personal honesty). Integrity = consistency of personal values.\\
Declining probity: Political (patronage), Administrative (poor work culture, weak accountability), Societal (normalization of corruption, materialism), Moral (declining values), Historical (bakshish system), Deterrence (poor punishment).\\
Philosophical basis: Utilitarianism, Deontology, Virtue Ethics, Weber, Social Contract, Rawls, Machiavelli, Gandhi, Kautilya, Bhagavad Gita.\\
Kautilya: Citizen-centric governance (`in their welfare, his welfare'); corruption is inevitable like fish drinking water; therefore create detection (spies) and deterrence (hefty punishment).
\end{QuickRevision}

\section{Utilization of Public Funds}

\subsection{Overview}
When government and bureaucrats utilize funds, those funds are public funds --- not from their personal pockets. Efficient and ethical utilization of these funds is a cornerstone of probity in governance. Problems with public fund utilization fall into two categories: under-utilization (funds not used effectively) and mis-utilization (funds diverted or misused).

\subsection{Reasons for Under-Utilization of Public Funds}

\paragraph{Procedural Bottlenecks and Red-Tapism.} Projects get stuck at various stages due to bureaucratic procedures --- environmental clearances not given, land acquisition delayed, approvals pending. When projects stall, allocated funds remain unspent or are spent without delivering outcomes.

\begin{ExampleBox}
\begin{itemize}
\item A bridge project was started over a seasonal river. The bridge was constructed but the approach road (land) was never acquired. So the bridge stands unused --- money was spent but the public cannot use it. This is under-utilization of funds due to lack of planning and procedural delays.
\item Cost escalation occurs when projects are delayed --- the original budget becomes insufficient, and additional funds are needed to complete the same project.
\end{itemize}
\end{ExampleBox}

\paragraph{Lack of Community Participation / Top-Down Implementation.} When schemes are designed at the center without consulting local communities, the expenditure may not be relevant to local needs.

\begin{ExampleBox}
\begin{itemize}
\item Centrally Sponsored Schemes (CSS): When the Planning Commission existed, states had to contribute a mandatory share (e.g., 40 crore out of 140 crore) even if they did not need that particular scheme. A state like Kerala might need more investment in tourism or land management rather than additional healthcare spending --- but the top-down approach forced committed expenditure on pre-determined schemes.
\item If the government builds many toilets in Karol Bagh (where they are not needed) instead of a library (which would be more useful), the funds are technically utilized but under-utilized in terms of optimal impact.
\end{itemize}
\end{ExampleBox}

\paragraph{Delay in Fund Disbursement.} When funds are not released on time, front-line workers (like ASHA workers) cannot plan activities. If funds arrive at the end of the financial year, there is insufficient time to implement programs effectively, leading to rushed and poor expenditure.

\paragraph{Incremental Budgeting.} Bureaucracy has a tendency to demand incrementally higher budgets each year without adequate justification. If a department received \rupee 100 crore last year, it will demand \rupee 110--120 crore this year --- simply because no bureaucrat asks for less. However, there may be no associated planning for the additional funds.

\begin{KeyConcept}
Incremental Budgeting: The practice of increasing budget demands year-on-year without justification. This leads to under-utilization because departments receive funds without having plans to use them effectively.\\
Zero-Based Budgeting: Recommended as an alternative --- every year's budget starts from zero, requiring fresh justification for every expenditure line item.
\end{KeyConcept}

\paragraph{Lack of Accountability.} Traditional accountability mechanisms like CAG operate in a post-mortem fashion and are institutional-centric. Day-to-day scrutiny of expenditure quality does not happen routinely, leading to wasteful spending going unchecked.

\subsection{Reasons for Mis-Utilization of Public Funds}

\paragraph{Corruption.} Bribery, embezzlement, and nepotism in tender processes all constitute mis-utilization. When tenders are given to friends or relatives (`bhaibanduon ko de diya'), the money goes to those it was never meant for. Corruption diverts public resources from their intended purpose.

\paragraph{Social Exclusion and Discriminatory Allocation.} When public funds are allocated without considering all sections of society, it constitutes mis-utilization. If \rupee 100 crore meant for public welfare is spent entirely on facilities for one gender or one region while excluding others, the fund is mis-utilized.

\begin{ExampleBox}
\begin{itemize}
\item If all public toilets built are male toilets, the social category of women is excluded. The \rupee 100 crore was never meant exclusively for male facilities.
\item Discriminatory allocation to politically significant states at the expense of others is also mis-utilization.
\end{itemize}
\end{ExampleBox}

\begin{GovtPolicy}
Gender-Based Budgeting: The practice of ensuring proportionate expenditure across genders. When making expenditure, government must ensure proportionate allocation for men, women, and third gender.
\end{GovtPolicy}

\paragraph{Ineffective Targeting of Beneficiaries.} When benefits reach the wrong people (ghost beneficiaries, fake entries), the purpose of the fund is defeated.

\begin{ExampleBox}
Ghost beneficiaries in MNAREGA --- names like Shahrukh Khan and Salman Khan in the muster rolls --- represent funds diverted through fake identities.
\end{ExampleBox}

\paragraph{Centralization of Decision-Making.} When decisions about fund utilization are made at the top without understanding local needs, resources may be directed to projects that are irrelevant to the community.

\subsection{Measures for Better Utilization}
\begin{itemize}
\item Zero-Based Budgeting: Justify every expenditure from scratch each year.
\item Outcome-Based Budgeting: Focus on outcomes, not just expenditure volumes.
\item Gender-Based Budgeting: Ensure proportionate allocation across demographics.
\item Strengthening social accountability mechanisms: RTI, social audit can improve transparency in fund utilization.
\item Community participation in planning: Bottom-up approach to expenditure planning.
\item Timely fund disbursement to implementing agencies.
\item Strengthening probity mechanisms: Anti-corruption institutions, transparent tender processes, technology-enabled monitoring.
\item Capacity building of implementing agencies.
\item Direct Benefit Transfer (DBT) to reduce leakages.
\end{itemize}

\begin{CriticalPoint}
Utilization of public funds connects directly to GS Paper III (Economics) topics like budgeting reforms. Cross-referencing between Paper III and Paper IV strengthens answers significantly.
\end{CriticalPoint}

\begin{QuickRevision}
Under-utilization: Procedural bottlenecks, lack of community participation, delayed fund disbursement, incremental budgeting, lack of accountability.\\
Mis-utilization: Corruption/nepotism, social exclusion/discriminatory allocation, ineffective targeting (ghost beneficiaries), centralized decision-making.\\
Measures: Zero-based budgeting, outcome-based budgeting, gender-based budgeting, social accountability mechanisms, community participation, timely disbursement, DBT, probity mechanisms.\\
Key insight: Incremental budgeting $\rightarrow$ Zero-based budgeting is a major reform recommendation.\\
Cross-reference with GS Paper III economics topics for stronger answers.
\end{QuickRevision}
